\uuid{Ojaf}
\exo7id{1806}
\titre{exo7 1806}
\auteur{gourio}
\organisation{exo7}
\datecreate{2001-09-01}
\isIndication{false}
\isCorrection{false}
\chapitre{Fonction de plusieurs variables}
\sousChapitre{Différentiabilité}
\cpge{mpsi-mp2i.fonctions-deux-variables.a, pcsi.fonctions-deux-variables.a, ptsi.fonctions-deux-variables.a, mpsi-mp2i.fonctions-deux-variables.b, pcsi.fonctions-deux-variables.b, ptsi.fonctions-deux-variables.b, bcpst1.fonctions-deux-variables.c, mp-mpi.topologie-espaces-normes.e, pc.espaces-vectoriels-normes.e, psi.espaces-vectoriels-normes.e, pt.fonctions-deux-trois-variables.b, mp-mpi.calcul-differentiel-optimisation.a}
\module{Analyse}
\niveau{L2}
\difficulte{}

\contenu{
\question{
Montrer qu'une norme $ N$ sur $ {\Rr}^{2} $ ne peut avoir des
d\'{e}riv\'{e}es partielles qui existent et qui soient continues en $ 0.$
}
}
